\scanpage{424}
\mspage{164}
aussitôt que l'on a comm. \unsure{du|des}
\diagram{\begin{tikzcd}[ampersand replacement=\&]
\widetilde{c'}\arrow[rr,"\widetilde\alpha"]\arrow[dr,"\widetilde{w'}"'] \& \& \widetilde{c'_{\unsure{1|i}}}\arrow[dl,"\widetilde{w'_1}"]\\
\&\widetilde c\& \& {\text{i.e.}\quad\widetilde{w'}=\widetilde{w'_1}\widetilde\alpha}
\end{tikzcd}}
i.e. que $\widetilde\alpha$ est bien une flèche de $\underline{\mathrm{Ch}}(X,\deleted{\unsure{x}}a',b')_{/\widetilde c}$. De $=$, que $\alpha\mapsto\widetilde\alpha$ définit une dépendance fonctorielle, i.e. $\widetilde{\beta\alpha}=\widetilde\beta\,\unsure{\widehat\alpha|\widetilde\alpha}$, est évident sur la définition de $\widetilde\alpha$. \deleted{\unsure{Cela}} Donc, on a bien défini
\[
\mathcal C(c,u,v)\xrightarrow{\varphi}\underline{\mathrm{Ch}}(X;a',b')_{/\widetilde c},
\]
définissons maintenant
\[
C(c,u,v)\xleftarrow{\psi}\underline{\mathrm{Ch}}(X;a',b')_{/\widetilde c}.
\]
Soit $c'',w''$ un objet de $\underline{\mathrm{Ch}}(X;a',b')_{/\widetilde c}$,
\diagram{\begin{tikzcd}[ampersand replacement=\&,row sep=small,column sep=1em]
a'_0\arrow[dr,equals]\&\&\&\makebox[0pt]{$\overbrace{\hspace{6.2cm}}^{c''}$}\&\&\\
\& a'\arrow[r]\arrow[d,dash,"\mathrm{id}"'] \& a'_1\arrow[rr,dash,"\rightarrow"]\arrow[d,"w''_1"] \& {}\arrow[d,dashed,"w''"] \& a'_n\arrow[d,"w''_n"'] \& b'\arrow[l]\arrow[d,equals]\\
\& a'\arrow[r,"u"] \& a\arrow[rr,dash] \& {} \& b \& b'\arrow[l,"v"'],\\
\&\&\&\makebox[0pt]{$\underbrace{\hspace{3.3cm}}_{c}$}\&\&
\end{tikzcd}}
donc
\[
w'':c''\mathrel{\unsure{\text{---}|\longrightarrow}}\widetilde c\qquad\text{flèche de }\underline{\mathrm{Ch}}(X;a',b');
\]
notons qu'on a une flèche de rétraction évidente
\[
\widetilde c\xrightarrow{\rho_c}c\deleted{\unsure{=\varphi}\,\unsure{\delta_{\Delta_+}}\unsure{\tau}\,\unsure{\delta_{\Delta_-}}}
\]
\inserted{associée à $\delta_{\tau}:\Delta_+\tau\Delta_-\to\tau$ (rétraction)} \unsure{induisant|traduisant} $u$ et $v$ sur la source et le but :
\diagram{\begin{tikzcd}[ampersand replacement=\&,column sep=.55em,row sep=small]
\deleted{\unsure{a}}\&a'\arrow[r,"{\overset{\frown}{u}}"]\arrow[d,"u"] \& a\arrow[rrr,dash,"c"]\arrow[d,dash,"\mathrm{id}"]\&{}\arrow[d,"\mathrm{id}"']\&{}\arrow[d,"\mathrm{id}"]\& b\arrow[d,"\mathrm{id}"]\&b'\arrow[l,"v"']\arrow[d,"v"]\&\big\}\ \widetilde c\in\mathrm{Ch}_{\Delta_+\tau\Delta_-}(a',b',\\
\&a\arrow[r,"\mathrm{id}"]\&a\arrow[rrr,dash,"c"]\&\&\&b\&b\arrow[l,"\mathrm{id}"']\&\big\}\ \delta_\tau^*(c)\in\mathrm{Ch}_{\Delta_+\tau\Delta_-}(a,b)\\
\&\&a\arrow[rrr,dash,"c","{\leftarrow}"']\&\&\&b\&\&\big\}\ c\in\mathrm{Ch}_\tau(a,b).
\end{tikzcd}}
% archive sheet number: 423
